% CONCLUSÃO--------------------------------------------------------------------

\chapter{CONCLUSÃO}
\label{chap:conclusao}

O presente trabalho resultou no desenvolvimento do aplicativo FishSpot, que teve como objetivo principal auxiliar os pescadores com dificuldades na localização de locais de pesca favoráveis, oferecendo uma solução prática e acessível. Este ambiente procura cumprir o objetivo de facilitar a busca por tais locais, demonstra potencial para impactar positivamente a comunidade pesqueira, fomentando a construção de laços entre seus usuários. 

Os recursos implementados incluem a possibilidade de encontrar registros de pesca em regiões específicas utilizando de um mapa digital, o registro de espécies de peixes e iscas utilizadas, a indicação da periculosidade dos locais e a visualização dos ultimos registros de pescarias realizadas.

Como trabalhos futuros, é possivel a expansão do aplicativo móvel para outras plataformas, visando ampliar o acesso e permitir que outros tipos de dispositivos consumam os recursos disponíveis no FishSpot. Adicionalmente, considera-se a implementação de uma rede social nativa no FishSpot. Tendo em vista que objetivo da aplicação é o compartilhamento de informações, este novo recurso poderia aumentar a distribuição de dados e a captação de novos usuários e parceiros, abrindo caminho para uma possível monetização. 

No entanto, para que estas futuras implementações sejam efetivadas, é crucial uma reformulação da interface atual para garantir a compatibilidade e a usabilidade com os novos recursos, bem como uma adaptação da arquitetura utilizada.
